% sections/11_2_factor_prices_draft.tex -- draft of subsection 11.2 (W2b) for W4 to \input into sections/11_scope.tex.
% Plan: BLUEPRINT 2.11 (11.2), 5.1-5.2; Decision on numbers: none reported here (no digits in the text; signs only).
% The algebra and the labeled equations eq:ge-prices, eq:ge-welfare, eq:ge-fiscal, eq:ge-mvpf and the sign table
% tab:gesigns are in appendix/D_derivations.tex (subsection app:factorprices); this text refers to them with \eqref/\autoref.
% W4: drop the \subsection line below if 11_scope.tex supplies its own heading; the label subsec:factorprices is referenced nowhere else yet.
\subsection{Factor prices: partial equilibrium and a reduced-form adjustment}\label{subsec:factorprices}

Every column of the reform tables is computed at the status-quo prices: the gross return $r_f$ and the age--health wage profiles $W(h,t)$ of \autoref{sec:model} do not respond to the cohort's choices. \autoref{subsec:scope} states the assumption as that of a small open economy, in which fixed $r$ and fixed wages are one assumption rather than two (Appendix~\ref{app:factorprices}). The reforms nevertheless change what the cohort supplies to factor markets: the participation columns of Tables~\ref{tab:claiming}--\ref{tab:other} report changes in labor supply by age group, and the saving subsidies change asset holdings (\autoref{subsec:saving}). In a closed economy, or if the reform applied to every cohort of a stationary economy, these changes would move the capital--labor ratio and with it the wage and the return. % src: T:welfare_savings_mvpf_v3.tex:120-123 (Indir. < 0 in every K row); 7:423 (taxes on asset income)

Appendix~\ref{app:factorprices} derives a reduced-form adjustment that requires no new solution of the model. With a constant-returns technology and competitive factor markets, the proportional change in the wage is the capital share divided by the elasticity of substitution, times the change in the capital--labor ratio, and the proportional change in the return gross of depreciation has the opposite sign and the labor share in place of the capital share, as in \eqref{eq:ge-prices}. At fixed behavior, the price changes alter each individual's earnings, net asset income, and, through the earnings record, retirement benefits. Valued at the relative marginal utilities of \autoref{sec:suffstats}, they add a term to the dollar willingness to pay, \eqref{eq:ge-welfare}, and, through the income tax, the payroll tax, and the benefit base, a term to the change in net spending, \eqref{eq:ge-fiscal}; the adjusted MVPF is the ratio of the adjusted numerator to the adjusted denominator, \eqref{eq:ge-mvpf}. The two terms partition the cohort's change in factor income between its members and the government.

The direction of the adjustment follows from the participation columns and is collected in \autoref{tab:gesigns}. Reforms that raise labor supply---the income-tax reductions at every age, and the claiming-schedule reforms after the normal retirement age---lower the wage and raise the return in a closed economy: workers lose part of the gain the tables record, asset holders gain, part of the indirect tax saving is returned, and the benefit base falls. Reforms that lower old-age labor supply---the larger benefit-formula increases and the Medicare expansions---raise the wage and lower the return, with the opposite incidence. The saving subsidies raise the capital stock and lower old-age participation, so they raise the wage and lower the return: the higher return they pay is partly undone, while labor-tax revenue and the benefit base rise. For the claiming-schedule reforms the sign of the aggregate labor-supply change depends on the weights, in efficiency units of labor (Appendix~\ref{app:factorprices}), of age groups whose participation moves in opposite directions. % src: signs of the participation columns in T:welfare_age_dependence_mvpfv5.tex:108-126, T:welfare_progressivity_mvpfv7.tex:115-158, T:welfare_medicare_mvpfv8.tex:113,116, T:welfare_tax_mvpf_v8.tex:114-118, T:welfare_savings_mvpf_v3.tex:120-123

No adjusted MVPF is reported. The participation changes in the tables are at most a few percentage points within age groups that are a minority of the cohort's labor in efficiency units, and the adjustment scales them by a factor share and divides by the substitution elasticity; the paragraph states the direction of the adjustment and not its size. A number would require the capital share, the elasticity of substitution, and the depreciation rate; the reform-induced changes in the cohort's labor in efficiency units and in its assets, computable from the simulated panels behind the tables; the age profiles of the labor-tax, asset-tax, and benefit bases in \eqref{eq:ge-fiscal}; and the elasticity of benefits to the wage profile. The companion paper lists factor-price feedback among the questions it leaves open \citep{ndiaye2026flexible}; the fixed-price environment used here is that of \citet{farhi2013insurance} and \citet{golosov2016redistribution}. % src: v6:9_conclusion_resubmission_v6.tex:10
